\textbf{Phần 1. Trắc nghiệm lựa chọn}
\setcounterpageref{ex}{0}

\begin{ex} %[5H9N50]
Công thức tính diện tích xung quanh của hình hộp chữ nhật có chiều dài $a$, chiều rộng $b$ và chiều cao $c$ (cùng đơn vị đo) là
\choice
{\True $S_{xq} = (a + b) \times 2 \times c$}
{$S_{xq} = a \times b \times c$}
{$S_{xq} = (a + b) \times c$}
{$S_{xq} = a \times b + c$}
\loigiai{
Diện tích xung quanh bằng chu vi đáy nhân với chiều cao: $S_{xq} = (a + b) \times 2 \times c$.\\
\textbf{Đáp án của câu hỏi là A}
}
\end{ex}

\begin{ex} %[5H9N50]
Diện tích toàn phần của hình hộp chữ nhật bằng
\choice
{Diện tích xung quanh cộng diện tích $1$ mặt đáy}
{\True Diện tích xung quanh cộng diện tích $2$ mặt đáy}
{Diện tích $4$ mặt bên}
{Chu vi đáy nhân chiều cao cộng chiều rộng}
\loigiai{
$S_{tp} = S_{xq} + S_{2đáy}$.\\
\textbf{Đáp án của câu hỏi là B}
}
\end{ex}

\begin{ex} %[5H9H50]
Một hình hộp chữ nhật có chiều dài $5\text{ cm}$, chiều rộng $3\text{ cm}$ và chiều cao $4\text{ cm}$. Diện tích xung quanh của hình hộp chữ nhật đó là
\choice
{$32\text{ cm}^2$}
{$60\text{ cm}^2$}
{$94\text{ cm}^2$}
{\True $64\text{ cm}^2$}
\loigiai{
$S_{xq} = (5 + 3) \times 2 \times 4 = 64\text{ (cm}^2\text{)}$.\\
\textbf{Đáp án của câu hỏi là D}
}
\end{ex}

\begin{ex} %[5H9H50]
Một hình hộp chữ nhật có chiều dài $8\text{ dm}$, chiều rộng $6\text{ dm}$ và chiều cao $5\text{ dm}$. Diện tích toàn phần của hình hộp chữ nhật đó là
\choice
{$140\text{ dm}^2$}
{$240\text{ dm}^2$}
{\True $236\text{ dm}^2$}
{$188\text{ dm}^2$}
\loigiai{
$S_{xq} = (8 + 6) \times 2 \times 5 = 140\text{ (dm}^2\text{)}$.\\
$S_{2đáy} = 2 \times 8 \times 6 = 96\text{ (dm}^2\text{)}$.\\
$S_{tp} = 140 + 96 = 236\text{ (dm}^2\text{)}$.\\
\textbf{Đáp án của câu hỏi là C}
}
\end{ex}

\begin{ex} %[5H9H50]
Một cái thùng tôn hình hộp chữ nhật không nắp có chiều dài $1{,}2\text{ m}$, chiều rộng $0{,}8\text{ m}$, chiều cao $0{,}6\text{ m}$. Diện tích tôn làm thùng (không kể mép hàn) là
\choice
{\True $3{,}36\text{ m}^2$}
{$2{,}4\text{ m}^2$}
{$4{,}32\text{ m}^2$}
{$2{,}88\text{ m}^2$}
\loigiai{
$S_{xq} = (1{,}2 + 0{,}8) \times 2 \times 0{,}6 = 2{,}4\text{ (m}^2\text{)}$.\\
$S_{đáy} = 1{,}2 \times 0{,}8 = 0{,}96\text{ (m}^2\text{)}$.\\
Diện tích tôn: $2{,}4 + 0{,}96 = 3{,}36\text{ (m}^2\text{)}$.\\
\textbf{Đáp án của câu hỏi là A}
}
\end{ex}

\begin{ex} %[5H9H50]
Một hình hộp chữ nhật có chu vi đáy là $24\text{ cm}$, chiều cao là $6\text{ cm}$. Diện tích xung quanh của hình hộp chữ nhật đó là
\choice
{$72\text{ cm}^2$}
{\True $144\text{ cm}^2$}
{$288\text{ cm}^2$}
{$120\text{ cm}^2$}
\loigiai{
$S_{xq} = P_{đáy} \times h = 24 \times 6 = 144\text{ (cm}^2\text{)}$.\\
\textbf{Đáp án của câu hỏi là B}
}
\end{ex}

\begin{ex} %[5H9V50]
Một phòng học hình hộp chữ nhật có chiều dài $8\text{ m}$, chiều rộng $6\text{ m}$ và chiều cao $3{,}5\text{ m}$. Người ta quét vôi tường xung quanh và trần nhà bên trong phòng học. Biết tổng diện tích các cửa sổ và cửa ra vào là $12\text{ m}^2$. Diện tích cần quét vôi là
\choice
{\True $134\text{ m}^2$}
{$146\text{ m}^2$}
{$144\text{ m}^2$}
{$182\text{ m}^2$}
\loigiai{
$S_{xq} = (8 + 6) \times 2 \times 3{,}5 = 98\text{ (m}^2\text{)}$.\\
$S_{trần} = 8 \times 6 = 48\text{ (m}^2\text{)}$.\\
Diện tích quét vôi: $98 + 48 - 12 = 134\text{ (m}^2\text{)}$.\\
\textbf{Đáp án của câu hỏi là A}
}
\end{ex}

\begin{ex} %[5H9V50]
Một hình hộp chữ nhật có diện tích xung quanh là $240\text{ cm}^2$, chiều cao $8\text{ cm}$, chiều dài hơn chiều rộng $3\text{ cm}$. Diện tích toàn phần của hình hộp chữ nhật đó là
\choice
{$300\text{ cm}^2$}
{$384\text{ cm}^2$}
{\True $348\text{ cm}^2$}
{$420\text{ cm}^2$}
\loigiai{
Nửa chu vi đáy: $240 : 8 : 2 = 15\text{ (cm)}$.\\
Chiều dài: $(15 + 3) : 2 = 9\text{ (cm)}$.\\
Chiều rộng: $15 - 9 = 6\text{ (cm)}$.\\
Diện tích 2 đáy: $2 \times 9 \times 6 = 108\text{ (cm}^2\text{)}$.\\
$S_{tp} = 240 + 108 = 348\text{ (cm}^2\text{)}$.\\
\textbf{Đáp án của câu hỏi là C}
}
\end{ex}

\begin{ex} %[5H9V50]
Người ta sơn toàn bộ các mặt ngoài của một chiếc hòm gỗ dạng hình hộp chữ nhật có nắp đậy với chiều dài $1{,}5\text{ m}$, chiều rộng $0{,}8\text{ m}$, chiều cao $0{,}6\text{ m}$. Chi phí sơn là $40\,000$ đồng/$\text{m}^2$. Tổng chi phí sơn chiếc hòm là
\choice
{$180\,000$ đồng}
{$220\,000$ đồng}
{$200\,000$ đồng}
{\True $206\,400$ đồng}
\loigiai{
$S_{xq} = (1{,}5 + 0{,}8) \times 2 \times 0{,}6 = 2{,}76\text{ (m}^2\text{)}$.\\
$S_{2đáy} = 2 \times 1{,}5 \times 0{,}8 = 2{,}4\text{ (m}^2\text{)}$.\\
$S_{tp} = 2{,}76 + 2{,}4 = 5{,}16\text{ (m}^2\text{)}$.\\
Chi phí: $5{,}16 \times 40\,000 = 206\,400$ (đồng).\\
\textbf{Đáp án của câu hỏi là D}
}
\end{ex}

\begin{ex} %[5H9V50]
Một bể bơi hình hộp chữ nhật có chiều dài $25\text{ m}$, chiều rộng $10\text{ m}$ và sâu $2\text{ m}$. Người ta lát gạch men xung quanh và đáy bể. Diện tích cần lát gạch là
\choice
{$390\text{ m}^2$}
{$400\text{ m}^2$}
{\True $390\text{ m}^2$}
{$500\text{ m}^2$}
\loigiai{
$S_{xq} = (25 + 10) \times 2 \times 2 = 140\text{ (m}^2\text{)}$.\\
$S_{đáy} = 25 \times 10 = 250\text{ (m}^2\text{)}$.\\
Diện tích lát gạch: $140 + 250 = 390\text{ (m}^2\text{)}$.\\
\textbf{Đáp án của câu hỏi là C}
}
\end{ex}

\begin{ex} %[5H9V50]
Nếu tăng cả chiều dài, chiều rộng và chiều cao của hình hộp chữ nhật lên gấp $2$ lần thì diện tích toàn phần của nó tăng lên gấp mấy lần?
\choice
{$2$ lần}
{\True $4$ lần}
{$6$ lần}
{$8$ lần}
\loigiai{
Khi các kích thước tăng 2 lần thì diện tích mỗi mặt tăng $2 \times 2 = 4$ lần, do đó diện tích toàn phần tăng lên 4 lần.\\
\textbf{Đáp án của câu hỏi là B}
}
\end{ex}

\begin{ex} %[5H9C50]
Một hình hộp chữ nhật có chiều cao $6\text{ cm}$, diện tích đáy bằng $45\text{ cm}^2$ và diện tích xung quanh bằng $168\text{ cm}^2$. Chiều dài của hình hộp chữ nhật đó là
\choice
{$7\text{ cm}$}
{$5\text{ cm}$}
{$8\text{ cm}$}
{\True $9\text{ cm}$}
\loigiai{
Nửa chu vi đáy: $168 : 6 : 2 = 14\text{ (cm)}$.\\
Ta có: $a + b = 14$ và $a \times b = 45$.\\
Hai số có tổng bằng $14$ và tích bằng $45$ là $9$ và $5$.\\
Vậy chiều dài là $9\text{ cm}$.\\
\textbf{Đáp án của câu hỏi là D}
}
\end{ex}

\textbf{Phần 2. Trắc nghiệm đúng sai}
\setcounterpageref{ex}{0}

\begin{ex} %[5H9H50]
Cho hình hộp chữ nhật có chiều dài $12\text{ cm}$, chiều rộng $8\text{ cm}$ và chiều cao $5\text{ cm}$. Xét tính đúng sai:
\choiceTF
{\True Chu vi mặt đáy của hình hộp là $40\text{ cm}$.}
{\True Diện tích xung quanh của hình hộp là $200\text{ cm}^2$.}
{Diện tích hai mặt đáy của hình hộp là $96\text{ cm}^2$.}
{\True Diện tích toàn phần của hình hộp là $392\text{ cm}^2$.}
\loigiai{
\begin{itemize}
    \item Chu vi đáy: $(12 + 8) \times 2 = 40\text{ cm}$. Đúng.
    \item $S_{xq} = 40 \times 5 = 200\text{ cm}^2$. Đúng.
    \item Diện tích 2 mặt đáy: $2 \times 12 \times 8 = 192\text{ cm}^2 \neq 96\text{ cm}^2$. Sai.
    \item $S_{tp} = 200 + 192 = 392\text{ cm}^2$. Đúng.
\end{itemize}
\begin{itemize}
    \item a) \textbf{Đúng}.
    \item b) \textbf{Đúng}.
    \item c) \textbf{Sai}.
    \item d) \textbf{Đúng}.
\end{itemize}
}
\end{ex}

\begin{ex} %[5H9V50]
Một thùng sắt không nắp hình hộp chữ nhật có chiều dài $2\text{ m}$, chiều rộng $1{,}5\text{ m}$, chiều cao $1\text{ m}$. Người ta sơn cả mặt trong lẫn mặt ngoài của thùng. Xét tính đúng sai:
\choiceTF
{\True Diện tích xung quanh của chiếc thùng là $7\text{ m}^2$.}
{\True Diện tích một mặt đáy của chiếc thùng là $3\text{ m}^2$.}
{Diện tích quét sơn một mặt (gồm xung quanh và 1 đáy) là $13\text{ m}^2$.}
{\True Tổng diện tích quét sơn cả trong và ngoài là $20\text{ m}^2$.}
\loigiai{
\begin{itemize}
    \item $S_{xq} = (2 + 1{,}5) \times 2 \times 1 = 7\text{ m}^2$. Đúng.
    \item $S_{đáy} = 2 \times 1{,}5 = 3\text{ m}^2$. Đúng.
    \item Diện tích 1 mặt sơn: $7 + 3 = 10\text{ m}^2 \neq 13\text{ m}^2$. Sai.
    \item Tổng diện tích sơn cả trong và ngoài: $10 \times 2 = 20\text{ m}^2$. Đúng.
\end{itemize}
\begin{itemize}
    \item a) \textbf{Đúng}.
    \item b) \textbf{Đúng}.
    \item c) \textbf{Sai}.
    \item d) \textbf{Đúng}.
\end{itemize}
}
\end{ex}

\textbf{Phần 3. Trả lời ngắn}
\setcounterpageref{ex}{0}

\begin{ex} %[5H9H50]
Một hình hộp chữ nhật có chiều dài $9\text{ cm}$, chiều rộng $7\text{ cm}$ và chiều cao $4\text{ cm}$. Tính diện tích xung quanh theo xăng-ti-mét vuông.
\shortans{128}
\loigiai{
\textbf{Đáp số:} $128$\\
$(9 + 7) \times 2 \times 4 = 128\text{ cm}^2$.
}
\end{ex}

\begin{ex} %[5H9H50]
Tính diện tích toàn phần của hình hộp chữ nhật có chiều dài $10\text{ dm}$, chiều rộng $5\text{ dm}$, chiều cao $6\text{ dm}$ theo đề-xi-mét vuông.
\shortans{280}
\loigiai{
\textbf{Đáp số:} $280$\\
$S_{xq} = (10 + 5) \times 2 \times 6 = 180\text{ dm}^2$.\\
$S_{2đáy} = 2 \times 10 \times 5 = 100\text{ dm}^2$.\\
$S_{tp} = 180 + 100 = 280\text{ dm}^2$.
}
\end{ex}

\begin{ex} %[5H9V50]
Người ta làm một cái hộp bằng bìa không nắp có chiều dài $15\text{ cm}$, chiều rộng $10\text{ cm}$, chiều cao $8\text{ cm}$. Tính diện tích bìa cần dùng theo xăng-ti-mét vuông.
\shortans{550}
\loigiai{
\textbf{Đáp số:} $550$\\
$S_{xq} = (15 + 10) \times 2 \times 8 = 400\text{ cm}^2$.\\
$S_{đáy} = 15 \times 10 = 150\text{ cm}^2$.\\
Diện tích bìa: $400 + 150 = 550\text{ cm}^2$.
}
\end{ex}

\begin{ex} %[5H9V50]
Một phòng học dài $7\text{ m}$, rộng $5\text{ m}$, cao $3{,}8\text{ m}$. Người ta quét vôi tường và trần nhà bên trong phòng học. Tổng diện tích các cửa sổ và cửa ra vào là $10{,}5\text{ m}^2$. Tính diện tích cần quét vôi theo mét vuông.
\shortans{115,7}
\loigiai{
\textbf{Đáp số:} $115{,}7$\\
$S_{xq} = (7 + 5) \times 2 \times 3{,}8 = 91{,}2\text{ m}^2$.\\
$S_{trần} = 7 \times 5 = 35\text{ m}^2$.\\
Diện tích quét vôi: $91{,}2 + 35 - 10{,}5 = 115{,}7\text{ m}^2$.
}
\end{ex}

\begin{ex} %[5H9V50]
Một hình hộp chữ nhật có diện tích xung quanh là $150\text{ cm}^2$, chiều cao là $5\text{ cm}$. Tính chu vi mặt đáy của hình hộp chữ nhật đó theo xăng-ti-mét.
\shortans{30}
\loigiai{
\textbf{Đáp số:} $30$\\
Chu vi đáy: $150 : 5 = 30\text{ cm}$.
}
\end{ex}

\begin{ex} %[5H9C50]
Một cái bể hình hộp chữ nhật không nắp có chu vi đáy $16\text{ m}$, chiều dài hơn chiều rộng $2\text{ m}$, chiều cao $1{,}5\text{ m}$. Người ta sơn các mặt bên ngoài của bể với giá $35\,000$ đồng/$\text{m}^2$. Tính tổng số tiền sơn bể theo đơn vị nghìn đồng.
\shortans{840}
\loigiai{
\textbf{Đáp số:} $840$\\
Nửa chu vi đáy: $16 : 2 = 8\text{ m}$. Chiều dài: $5\text{ m}$, chiều rộng: $3\text{ m}$.\\
$S_{xq} = 16 \times 1{,}5 = 24\text{ m}^2$.\\
Tiền sơn ngoài xung quanh: $24 \times 35 = 840$ nghìn đồng.
}
\end{ex}
